\section{Introduction}

Large language models (LLMs) are becoming physical operators in data systems.
Systems such as LOTUS~\cite{patel2025lotus},
Palimpzest~\cite{liu2025palimpzest}, UQE~\cite{dai2024uqe}, and
DocETL~\cite{shankar2025docetl} expose declarative operators over text, images,
and documents, then choose an execution strategy that trades financial cost,
latency, and result quality.
For a bulk semantic filter or map, even a modest per-tuple price difference is
multiplied by the cardinality of the input relation. Always using the strongest
available model preserves quality but can dominate query cost; always using a
cheap model is efficient but can corrupt downstream filters, joins, and
aggregates.

Model routing treats the model choice as an execution-plan decision. Given a
tuple and operator, a router selects either a cheap model or a stronger model.
Prior LLM-routing work learns this decision from confidence, preferences, or
model features and reports favorable aggregate cost--quality curves
\cite{chen2024frugalgpt,ong2025routellm,gupta2024cascades,jitkrittum2026uniroute}.
Bulk semantic processing creates a particularly unforgiving version of the
problem. The router is invoked repeatedly, failures may be correlated within a
group or predicate, and a small number of false negatives can materially alter
an aggregate. Moreover, the useful training signal is severely imbalanced:
the cheap and strong model agree on most tuples, while the operationally
important event---cheap wrong and strong correct---may be rare.

Consider a sentiment operator over movie reviews. For a review saying
``the film is a disaster, except that it falls into the so-bad-it-is-good
category,'' both models may agree, only the strong model may resolve the
contrast, only the cheap model may happen to be right, or both may fail.
Difficulty alone does not determine the desirable route. A difficult tuple on
which both models fail cannot gain answer accuracy from escalation; conversely,
an apparently simple tuple can be a critical cheap-model failure. A router
trained only on easy/hard labels, overall route accuracy, or class balance
therefore optimizes a surrogate that is misaligned with the query result.

This paper focuses on \emph{lightweight soft-prompt routing} as the primary
way to learn this execution decision. A frozen compact encoder maps each tuple
to a routing representation; only a short continuous prompt and classification
head are trained. This avoids fine-tuning the answer models and, unlike an
LLM-based router, avoids autoregressive decoding by a multi-billion-parameter
model on every tuple. We study discrete natural-language prompt optimization
separately as an interpretable alternative, not as a stage of the soft-prompt
router. Existing prompt optimizers show substantial task gains
\cite{pryzant2023protegi,opsahlong2024mipro,agrawal2025gepa}, but their usual
training data and scalar metrics do not target rare recoverable failures or
distinguish a harmless expensive call from a critical under-escalation.

We introduce \system around two methodological innovations. First,
\emph{disagreement-guided critical-case acquisition} starts from an unlabeled
candidate pool and executes both answer models without revealing gold. It
partitions candidates by whether the two predictions agree, prioritizes
disagreements, retains an agreement sample for coverage, and requests human
gold labels only for the selected subset. Every Mini-only and Nano-only case
must appear in the disagreement stratum, so this observable signal concentrates
annotation on routing-critical examples without inspecting hidden correctness.
On Computer Science QA, both disagreement-guided and random acquisition screen
the same 200 pairs and request 100 human labels. Disagreement-guided selection
covers all 19 Mini-only and all 8 Nano-only cases; random selection covers
53.1\% and 54.8\% on average. Thus the saving is 50\% of human labels relative
to labeling the full screened pool, not a saving in paired model calls. Second,
a \emph{cost-derived route score} assigns inverse expected route cost to a
correct route and zero to an incorrect route, making every routed decision
auditable in common units. The primary soft-prompt router is trained only with
inverse-expected-cost weighted cross-entropy. Disagreement enrichment, rather
than a hand-tuned Mini-only weight, ensures that rare recoverable errors
influence training. The cost score reports the resulting cost--quality frontier
and supplies feedback to the separate GEPA adapter; it is not a second
soft-prompt loss.
The empirical finding is deliberately more nuanced than ``routing saves
cost.'' The lightweight soft-prompt router has only 3,074--13,826 trainable
parameters and provides a smooth validation cost--quality curve. On historical
Movie splits, outcome enrichment and cost-weighted CE produce configurations
near all-Mini accuracy while avoiding 15.85--34.15\% of Mini calls. However, the
locked fresh evaluation shows that the apparent frontier does not robustly
reproduce: the mean router loses 1.33 accuracy points and Mini-call savings vary
by 24 points across seeds. More soft-prompt tokens do not monotonically help.
The independent discrete router has a different failure mode: hard routing
collapses to all-Nano, and a calibrated verbal risk score only partly repairs
it. Computer Science confirms that preserving all-Mini accuracy remains
expensive under the present data budget.

Our methodological novelty is concentrated in the first two contributions:

\begin{itemize}[leftmargin=*]
  \item \textbf{Disagreement-guided acquisition of critical routing data.}
  We screen an unlabeled pool with both answer models, prioritize candidates
  whose answers disagree, retain agreement coverage, and obtain human gold
  labels only for the selected subset. The observable split enriches the
  one-model-only cases that control router accuracy and reduces annotation
  demand without using correctness during selection.
  \item \textbf{Weighted soft-prompt routing with cost-derived evaluation.}
  Paired execution assigns both routes an auditable score from correctness and
  expected route cost. Inverse-expected-cost weighted cross-entropy trains only
  3,074--13,826 parameters over a frozen compact encoder. The score also
  supplies feedback through a separate GEPA adapter.

  \item \textbf{Independent comparison and careful evaluation.} Across 1,012 paired Movie reviews, a locked 200-row fresh set,
  and 410 Computer Science questions, we compare the lightweight soft-prompt
  method against reflective discrete instructions, loss designs, acquisition
  policies, and validation thresholds.
\end{itemize}

The current draft is an Experiment, Analysis \& Benchmark submission because
its strongest result is the controlled characterization of reliability, not a
claim that one prototype dominates all baselines. The final artifact will
release split manifests, paired outcomes, optimization traces, and scripts
needed to reproduce every table.
